\begin{comment}
  Homework solution of Calculus 2 - Chapter 4 Integral Application
  Licensed under MIT license
\end{comment}

\documentclass[12pt]{article}
\usepackage[utf8]{inputenc}
\usepackage{amsmath}
\usepackage{geometry}

\geometry{a4paper, margin=1in}

\title{No. 7 - Aplikasi Integral}
\author{Ahmad Fakhrul Bawani}
\date{}

\begin{document}

\maketitle

\section*{Solusi Volume Benda Pejal}

Benda pejal terletak di antara $y = 0$ dan $y = 1$. Penampang melintang yang tegak lurus sumbu-$y$ berbentuk piringan lingkaran (\textit{circular disk}) dengan diameter yang membentang dari sumbu-$y$ ($x = 0$) hingga kurva parabola $x = \sqrt{14}y^2$.

\subsection*{1. Menentukan Jari-jari Lingkaran ($r$)}
Diameter penampang ($D$) adalah jarak horizontal dari sumbu-$y$ ke kurva:
\[ D(y) = \sqrt{14}y^2 - 0 = \sqrt{14}y^2 \]
Jari-jari ($r$) adalah setengah dari diameter:
\[ r(y) = \frac{\sqrt{14}y^2}{2} \]

\subsection*{2. Menentukan Luas Penampang ($A$)}
Karena penampang berbentuk lingkaran, luasnya ($A$) adalah:
\[ A(y) = \pi [r(y)]^2 = \pi \left( \frac{\sqrt{14}y^2}{2} \right)^2 \]
\[ A(y) = \pi \left( \frac{14y^4}{4} \right) = \frac{7\pi}{2} y^4 \]

\subsection*{3. Menghitung Volume ($V$)}
Volume diperoleh dengan mengintegralkan fungsi luas penampang terhadap $y$ dari $0$ hingga $1$:
\begin{align*}
V &= \int_{0}^{1} A(y) \, dy \\
V &= \int_{0}^{1} \frac{7\pi}{2} y^4 \, dy \\
V &= \frac{7\pi}{2} \left[ \frac{1}{5} y^5 \right]_{0}^{1} \\
V &= \frac{7\pi}{2} \left( \frac{1}{5}(1)^5 - 0 \right) \\
V &= \frac{7\pi}{10}
\end{align*}

\subsection*{Kesimpulan}
Jadi, volume benda pejal tersebut adalah $\mathbf{\frac{7\pi}{10}}$ \textbf{satuan kubik}.

\end{document}